\section{Two raw realizations without a positivity floor}
\label{sec:raw}
\subsection*{Paid, history-dependent heralding}
There are two physically distinct ports. At port $G$ a gate call costs one. Its success probability, conditional on every hidden and observed past, belongs to $[p_0,1]$, where $p_0>0$ is a known guaranteed lower bound, not the unknown true rate. Success permits one call at acquisition port $A$, costing one. No commitment is legal at $G$. After rejecting at $A$, the protocol returns to $G$. No hidden state is reset by this routing convention. The class includes a reference model with constant gate success $p_0$ and the productive laws specified below.

\begin{lemma}[The exact robust paid-attempt scale]\label{lem:gate}
For this interface set $\kappa=1+1/p_0$. Every all-history policy under the reference model satisfies $\E C=\kappa\E A_\tau$. The program which calls the legal current port reaches one productive report at conditional expected cost at most $\kappa$, for every model and history. Both ports require distinct live states, so $B=2$ is a semantic mandatory lower. Robust expected-budget feasibility is exactly $N\ge\kappa$ when at least three observer states are available.
\end{lemma}
\begin{proof}
Let $G_\tau,A_\tau$ be gate and productive counts. Every productive call consumes one gate success, and a terminating path has no unused success. Predictable counting gives $\E A_\tau=p_0\E G_\tau$ under the witness, without requiring independent cycles. Hence $\E C=(1+1/p_0)\E A_\tau$. For the upper, the probability of $j$ consecutive failures, conditional on any history at the start, is at most $(1-p_0)^j$. Summing gives at most $1/p_0$ gate calls and one acquisition call. The physical gate and acquisition ports have disjoint legal command sets, and both occur before any legal commitment under the same witness. Proposition~\ref{prop:physical} applies. Accepting the first productive report gives the feasibility converse. At $p_0=1$ the cost is two, not one. At fixed $N$, the task is infeasible when $1+1/p_0>N$; rare-rate uniformity always refers to the feasible rescaled budget $N/\kappa$.
\end{proof}

\subsection*{Controlled classical hidden states}
Let the physical state contain a sign $E\in\{-1,1\}$ and an arbitrary nuisance state. Between productive calls any causal stochastic or deterministic evolution is allowed, with no mixing or positivity requirement. The observer may choose either of two physical commands on the sign: identity or sign flip, followed immediately by the same known observation-and-audit instrument. There is no unknown intervening evolution between this chosen command and that instrument. The physical task is to predict the \emph{post-command} audit, not a pre-intervention sign.

For input sign $e$, let $J_e(dx,dz)$ be any fixed probability kernel with bounded audit in $H$; it can have atoms, singular parts and sign-dependent report likelihoods. Nuisance evolution is unrestricted provided this productive instrument depends on the sign through these stated kernels. The class contains a reference member which supplies a fresh fair sign before each productive command, independently of all prior reports, and uses the constant gate rate $p_0$. This is an admitted physical preparation model, not a reset available freely to every observer. Its productive law is
\begin{equation}\label{eq:classical-ref}
 K_0=\tfrac12(J_++J_-).
\end{equation}

\begin{theorem}[Classical task balance at zero positivity floor]\label{thm:classical}
The classical family satisfies Theorem~\ref{thm:main} with $B=2$, $D_0=1$ and $\kappa=1+1/p_0$, for the task geometry derived from \eqref{eq:classical-ref}. In particular, for all $N\ge\kappa$ and $M\ge3$,
\begin{equation}\label{eq:classical-risk}
 R_\aut^\Theta(N,M)=\Qm_{M-2}(N/\kappa),\qquad
 R_\cp^\Theta(N,m)=\Qm_m(N/\kappa).
\end{equation}
The constants require no positive lower bound on transition probabilities or conditional residual variances.
\end{theorem}
\begin{proof}
The upper chooses identity or flip with a fresh fair coin within the productive command. For every conditional pre-command sign distribution, averaging makes the post-command sign fair. By linearity the joint report--audit law, not merely its report marginal, is exactly $K_0$. Take the identity export map. The nuisance state need not be reset or stored. Under the reference member, each chosen physical command, even conditioned on its chosen coin and on the full preceding history, still has law $K_0$; hence arbitrary competitors satisfy the reference conditions. Lemma~\ref{lem:gate} gives both cost directions and the mandatory states. Theorem~\ref{thm:main} proves \eqref{eq:classical-risk}. Classical instrument continuity follows from Proposition~\ref{prop:instrument}, including deterministic transitions.
\end{proof}

For a concrete singular geometry, draw $(Y,T)$ independently of the sign with any probability $\nu_0$ on $K\times[0,1]$, and let $Z=Y+TEe_1$ in a fixed compact convex carrier. The reference has $f=Y,w=T^2$. Choose a $d$-regular singular measure on $K$ and reliability law $\Pp(T\le s)=s^\alpha$ to obtain \eqref{eq:rectangle}; add a genuinely prepared atom $(Y,T)=(x_*,0)$ to obtain Proposition~\ref{prop:atom}. The report can instead be informative: a weak sign report $V$ with $\Pp(V=v\mid E=e)=(1+\eta ve)/2$ gives $f=Y+\eta VT e_1$ and $w=(1-\eta^2)T^2$ under the reference. Its profile is still exact, although rectangle estimates must then be checked for that sheared marked law rather than assumed. In contrast, a static known-to-the-history sign can have zero conditional variance at individual histories. Such histories violate a positive pointwise variance mark but do not invalidate this minimax theorem.

\subsection*{Noncommuting quantum instruments with rank changes}
Consider a qubit with arbitrary unknown completely positive trace-preserving evolution between port interactions. No depolarizing component, spectral gap or full-rank floor is imposed. Gate instruments have success effects bounded below by $p_0I$; their surviving states may be pure. The productive port permits a Pauli command from $\{I,\sigma_x,\sigma_y,\sigma_z\}$, followed by a fixed known instrument and audit, with no unknown operation inserted inside this command. The class contains a reference member that prepares $I/2$ immediately before every productive Pauli command and has gate success $p_0$.

Here is an explicit noncommuting instrument. Draw $(y,t)$ with any fixed probability $\nu_0$ on $K\times[0,1]$. Choose a Borel angle $\phi(y,t)$ and $0\le\eta(y,t)\le1$. Set
\[
 n=(\cos\phi,\sin\phi,0),\qquad
 E_v=\tfrac12(I+v\eta n\cdot\sigma),\quad L_v=\sqrt{E_v},\qquad v=\pm1.
\]
The unnormalized post-report state is $L_v\rho L_v$, and the report is $x=(y,t,v)$. At commitment measure $\sigma_z$, with result $e=\pm1$, and score $Z=y+te e_1$. The axes may vary, so effects at different reports need not commute.

\begin{theorem}[Quantum verification through pure posterior strata]\label{thm:quantum}
For the stated quantum family, the same conclusions \eqref{eq:classical-risk} hold with the actual reference mark law
$\nu=\nu_0\otimes\tfrac12(\delta_{-1}+\delta_1)$ and $f=y,w=t^2$. They remain valid at $\eta=1$, at pure incoming states, and across report strata where the posterior rank changes. If $\nu_0$ has the geometry or atom described above, the corresponding regular or exact atomic risk law follows from the \emph{same} central theorem.
\end{theorem}
\begin{proof}
Writing $\rho=(I+r\cdot\sigma)/2$ shows directly that
\[
 \tfrac14\sum_{U\in\{I,\sigma_x,\sigma_y,\sigma_z\}}U\rho U^*=I/2.
\]
This standard Pauli averaging identity \cite{AMTW} is a charged within-call preparation, not free additional quantum memory. Every unnormalized report--audit probability is linear in $\rho$, so its average is its value at $I/2$. The reference member is invariant under \emph{each} chosen Pauli command, including after conditioning on that command and all earlier observations. Thus it gives the required reference for all competitors, while the random command gives one common upper for every unknown input state.

Under $I/2$, the probability of $v$ is $\tr(E_v)/2=1/2$ and the normalized post-report state is $E_v$. Its $z$ component is zero, so the subsequent Pauli audit has both signs with probability $1/2$. This gives the stated $f,w$ and the full unnormalized loss identity. The eigenvalues of $E_v$ are $(1\pm\eta)/2$. Choosing $\eta=1$ on one set of positive $\nu_0$-mass and $\eta=1/2$ on another gives actual rank-one and rank-two strata, not merely a formal rank-deficient limit. Different equatorial axes give noncommuting effects. Proposition~\ref{prop:instrument} supplies report continuity and actual-mass-weighted update stability without dividing by a positive lower eigenvalue. Lemma~\ref{lem:gate} and Theorem~\ref{thm:main} complete the proof.
\end{proof}

The quantum carrier here is finite dimensional and the observer is classical. It has two live labels and its terminal dictionary, not an uncharged quantum register. Interventions change the declared post-command audit; neither theorem asserts that randomizing an unknown pre-intervention quantity preserves its original estimation task. Known classical or quantum preparation operations, their costs and the availability of the reference member are explicit hypotheses, not automatic consequences of positivity.
